% GENERATED FILE. Edit canonical lessons, then run npm run generate:books.
% canonical-source-hash: fnv1a64:d9d60e19bc5382f6
% canonical-lessons: UR-C498-chak, UR-C498-sancha, UR-C498-ava, UR-C498-matka, UR-C498-tapana

\chapter{Clay and the Kiln}
\label{ch:ur-clay-and-the-kiln}
\hlchaptermodality{498}

\begin{chapteropening}
\textbf{By the end of this chapter:} \emph{I can describe how a potter shapes and fires pots.}

Complete the last lesson of chapter 498: I can describe how a potter shapes and fires pots.
\end{chapteropening}

\section[\texorpdfstring{\ur{چاک} (chāk)}{چاک (chāk)}]{\ur{چاک} (chāk) — a potter's wheel}

\label{lesson:UR-C498-chak}



\begin{quote}
Before the new one: say the Urdu for to knock in, then the Urdu for to file (metal or wood).
\end{quote}

\subsection*{You'll want to know: \ur{چاک}}
\textbf{\ur{چاک}} — \emph{chāk} — \textquotedblleft{}a potter's wheel\textquotedblright{}.

\begin{quote}
\ur{کمہار} \ur{چاک} \ur{پر} \ur{مٹی} \ur{کا} \ur{برتن} \ur{بناتا} \ur{ہے۔}
\end{quote}

\begin{quote}
\emph{kumhār chāk par miṭṭī kā bartan banātā hai.}
\end{quote}

\begin{quote}
The potter makes a clay pot on the wheel.
\end{quote}

Masculine.

\begin{cousinweb}
\textbf{\ur{چاک}} (\emph{chāk}) is inherited: CDIAL lists the Hindi cousin \emph{cāk}, \textquotedblleft{}any kind of wheel\textquotedblright{}, under Sanskrit \emph{cakra}, \textquotedblleft{}a wheel\textquotedblright{} (CDIAL 4538).
\end{cousinweb}

\begin{grammarlens}[title={What you've built}]
One more word for describing what you lived and felt.
\end{grammarlens}

\subsection*{Your turn}
\begin{itemize}
\raggedright
  \item \emph{Say it:} \emph{chāk}
  \item \emph{Say it:} \emph{chāk}, once more
  \item \emph{Say it:} \emph{kumhār chāk par miṭṭī kā bartan banātā hai.} — The potter makes a clay pot on the wheel.
  \item \emph{Recall:} say the Urdu for to knock in, then the Urdu for to file (metal or wood), then say \emph{chāk} again
\end{itemize}

\begin{tcolorbox}[breakable,colback=teal!4,colframe=teal!35!black,title={Before you move on}]
What does \ur{چاک} mean? (\textquotedblleft{}A potter's wheel\textquotedblright{}.) Say it once more.
\end{tcolorbox}
\section[\texorpdfstring{\ur{سانچا} (sāṅchā)}{سانچا (sāṅchā)}]{\ur{سانچا} (sāṅchā) — a mould}

\label{lesson:UR-C498-sancha}



\begin{quote}
Before the new one: say the Urdu for to file (metal or wood), then the Urdu for a potter's wheel.
\end{quote}

\subsection*{You'll want to know: \ur{سانچا}}
\textbf{\ur{سانچا}} — \emph{sāṅchā} — \textquotedblleft{}a mould\textquotedblright{}.

\begin{quote}
\ur{اس} \ur{سانچے} \ur{سے} \ur{مٹی} \ur{کے} \ur{کھلونے} \ur{بنتے} \ur{ہیں۔}
\end{quote}

\begin{quote}
\emph{is sāṅche se miṭṭī ke khilaune bante haiṅ.}
\end{quote}

\begin{quote}
Clay toys are made in this mould.
\end{quote}

Masculine.

\begin{cousinweb}
\textbf{\ur{سانچا}} (\emph{sāṅchā}) is inherited: CDIAL lists the Hindi cousin \emph{sā̃cā}, \textquotedblleft{}a mould\textquotedblright{}, under Sanskrit \emph{sañcaka}, \textquotedblleft{}a mould, a figure\textquotedblright{} (CDIAL 13096).
\end{cousinweb}

\begin{grammarlens}[title={What you've built}]
One more word for describing what you lived and felt.
\end{grammarlens}

\subsection*{Your turn}
\begin{itemize}
\raggedright
  \item \emph{Say it:} \emph{sāṅchā}
  \item \emph{Say it:} \emph{sāṅchā}, once more
  \item \emph{Say it:} \emph{is sāṅche se miṭṭī ke khilaune bante haiṅ.} — Clay toys are made in this mould.
  \item \emph{Recall:} say the Urdu for to file (metal or wood), then the Urdu for a potter's wheel, then say \emph{sāṅchā} again
\end{itemize}

\begin{tcolorbox}[breakable,colback=teal!4,colframe=teal!35!black,title={Before you move on}]
What does \ur{سانچا} mean? (\textquotedblleft{}A mould\textquotedblright{}.) Say it once more.
\end{tcolorbox}
\section[\texorpdfstring{\ur{آوا} (āvā)}{آوا (āvā)}]{\ur{آوا} (āvā) — a potter's kiln}

\label{lesson:UR-C498-ava}



\begin{quote}
Before the new one: say the Urdu for a potter's wheel, then the Urdu for a mould.
\end{quote}

\subsection*{You'll want to know: \ur{آوا}}
\textbf{\ur{آوا}} — \emph{āvā} — \textquotedblleft{}a potter's kiln\textquotedblright{}.

\begin{quote}
\ur{کمہار} \ur{نے} \ur{برتن} \ur{آوے} \ur{میں} \ur{رکھے۔}
\end{quote}

\begin{quote}
\emph{kumhār ne bartan āve meṅ rakhe.}
\end{quote}

\begin{quote}
The potter put the pots in the kiln.
\end{quote}

Masculine.

\begin{cousinweb}
\textbf{\ur{آوا}} (\emph{āvā}) is inherited from Sanskrit \emph{āpāka}, \textquotedblleft{}a baking oven, a potter's kiln\textquotedblright{} (CDIAL 1202).
\end{cousinweb}

\begin{grammarlens}[title={What you've built}]
One more word for describing what you lived and felt.
\end{grammarlens}

\subsection*{Your turn}
\begin{itemize}
\raggedright
  \item \emph{Say it:} \emph{āvā}
  \item \emph{Say it:} \emph{āvā}, once more
  \item \emph{Say it:} \emph{kumhār ne bartan āve meṅ rakhe.} — The potter put the pots in the kiln.
  \item \emph{Recall:} say the Urdu for a potter's wheel, then the Urdu for a mould, then say \emph{āvā} again
\end{itemize}

\begin{tcolorbox}[breakable,colback=teal!4,colframe=teal!35!black,title={Before you move on}]
What does \ur{آوا} mean? (\textquotedblleft{}A potter's kiln\textquotedblright{}.) Say it once more.
\end{tcolorbox}
\section[\texorpdfstring{\ur{مٹکا} (maṭkā)}{مٹکا (maṭkā)}]{\ur{مٹکا} (maṭkā) — a large clay pot}

\label{lesson:UR-C498-matka}



\begin{quote}
Before the new one: say the Urdu for a mould, then the Urdu for a potter's kiln.
\end{quote}

\subsection*{You'll want to know: \ur{مٹکا}}
\textbf{\ur{مٹکا}} — \emph{maṭkā} — \textquotedblleft{}a large clay pot\textquotedblright{}.

\begin{quote}
\ur{گرمیوں} \ur{میں} \ur{ہم} \ur{مٹکے} \ur{میں} \ur{پانی} \ur{رکھتے} \ur{ہیں۔}
\end{quote}

\begin{quote}
\emph{garmiyoṅ meṅ ham maṭke meṅ pānī rakhte haiṅ.}
\end{quote}

\begin{quote}
In summer we keep water in a big clay pot.
\end{quote}

Masculine.

\begin{cousinweb}
\textbf{\ur{مٹکا}} (\emph{maṭkā}) is inherited: CDIAL lists the Hindi cousin \emph{maṭkā}, \textquotedblleft{}an earthen jar\textquotedblright{}, under Sanskrit \emph{mṛdā}, \textquotedblleft{}earth, clay\textquotedblright{} (CDIAL 10291).
\end{cousinweb}

\begin{grammarlens}[title={What you've built}]
One more word for describing what you lived and felt.
\end{grammarlens}

\subsection*{Your turn}
\begin{itemize}
\raggedright
  \item \emph{Say it:} \emph{maṭkā}
  \item \emph{Say it:} \emph{maṭkā}, once more
  \item \emph{Say it:} \emph{garmiyoṅ meṅ ham maṭke meṅ pānī rakhte haiṅ.} — In summer we keep water in a big clay pot.
  \item \emph{Recall:} say the Urdu for a mould, then the Urdu for a potter's kiln, then say \emph{maṭkā} again
\end{itemize}

\begin{tcolorbox}[breakable,colback=teal!4,colframe=teal!35!black,title={Before you move on}]
What does \ur{مٹکا} mean? (\textquotedblleft{}A large clay pot\textquotedblright{}.) Say it once more.
\end{tcolorbox}
\section[\texorpdfstring{\ur{تپانا} (tapānā)}{تپانا (tapānā)}]{\ur{تپانا} (tapānā) — to heat; to fire (pots)}

\label{lesson:UR-C498-tapana}



\begin{quote}
Before the new one: say the Urdu for a potter's kiln, then the Urdu for a large clay pot.
\end{quote}

\subsection*{You'll want to know: \ur{تپانا}}
\textbf{\ur{تپانا}} — \emph{tapānā} — \textquotedblleft{}to heat; to fire (pots)\textquotedblright{}.

\begin{quote}
\ur{کمہار} \ur{برتنوں} \ur{کو} \ur{آگ} \ur{میں} \ur{تپاتا} \ur{ہے۔}
\end{quote}

\begin{quote}
\emph{kumhār bartanoṅ ko āg meṅ tapātā hai.}
\end{quote}

\begin{quote}
The potter fires his pots in the flames.
\end{quote}

Past, with \ur{نے} (he or she …): \ur{اس} \ur{نے} \ur{برتن} \ur{آگ} \ur{میں} \ur{تپائے}, \emph{us ne bartan āg meṅ tapāe}.

\begin{cousinweb}
\textbf{\ur{تپانا}} (\emph{tapānā}) is the causative of \emph{tapnā}, \textquotedblleft{}to be hot\textquotedblright{}, which goes back to Sanskrit \emph{tapyati}, \textquotedblleft{}is hot\textquotedblright{} (CDIAL 5684).
\end{cousinweb}

\begin{grammarlens}[title={What you've built}]
One more word for describing what you lived and felt.
\end{grammarlens}

\subsection*{Your turn}
\begin{itemize}
\raggedright
  \item \emph{Say it:} \emph{tapānā}
  \item \emph{Say it:} \emph{tapānā}, once more
  \item \emph{Say it:} \emph{kumhār bartanoṅ ko āg meṅ tapātā hai.} — The potter fires his pots in the flames.
  \item \emph{Recall:} say the Urdu for a potter's kiln, then the Urdu for a large clay pot, then say \emph{tapānā} again
\end{itemize}

\begin{tcolorbox}[breakable,colback=teal!4,colframe=teal!35!black,title={Before you move on}]
What does \ur{تپانا} mean? (\textquotedblleft{}To heat; to fire (pots)\textquotedblright{}.) Say it once more.
\end{tcolorbox}
